\section{Introduction}
\label{sec:intro}
Single-writer value transfer does not require every payment to occupy a position in a global transaction sequence. An account owner determines the order of its updates; validators must prevent incompatible owner-signed histories from becoming final. The consensus-number characterization of asset transfer formalizes this distinction, and payment systems already exploit it~\cite{consensusnumber,fastpay,astro}. The challenge addressed here is not ordinary consensus avoidance. It is replacing the validators that have been issuing account-local votes.

A reconfiguration must determine which old outcomes the successor is obliged to preserve. Consider a payment for which enough old validators have released signatures to establish finality. Network delays keep these signatures distributed, so a checkpoint builder sees no final certificate. If the checkpoint omits the payment and allows an incompatible successor history, the delayed signatures can subsequently establish a valid old-epoch final outcome. Agreement has selected a unique snapshot, but not a safe one.

We call an eligible outcome \emph{latent-final} when the released signatures suffice for a final certificate, whether or not a correct process has assembled it. More generally, closure must preserve outcomes that can still become certified using signatures from validators that have not yet stopped. An immutable report quorum must expose enough obligations to cover both cases. The relevant boundary is therefore not ``which certificates did this proposer observe?'' but ``which outcomes can the old voting domain still justify?''

RAI separates \emph{closure}, which recovers these obligations, from \emph{checkpoint agreement}, which chooses among safe candidates. A closing validator stops account signing and freezes one report containing two authenticated roots. The first commits to its ledger with distinct recovery-protected, notarized, and finalized tags. The second commits to hashes of blocks for which it voted but which are absent from that frozen ledger. Reports do not contain certificate objects. Validators retain and gossip blocks and signed votes, from which every verifier assembles quorum evidence locally.

A candidate selects $N-f$ usable reports. The deterministic procedure $\Build$ preserves inherited and explicitly exposed finality, carries represented notarization locks, and recovers a potentially fast-final branch from sufficient reporter first votes. After mandatory retention, a decided checkpoint may additionally finalize the maximal uniquely retained prefix whose blocks each have a verified closing-epoch $\NC$. Recovery-only uniqueness remains nonfinal. The crucial property is universal over valid report selections: different candidates may retain different nonfinal work, but none may discard a branch on which eligible late finality can act.

This design specializes entirely to owner-ordered sends and receives. It does not provide arbitrary shared-object transactions or atomic operations across independent writers. Sui Lutris, in contrast, combines owned-object execution with consensus for shared objects and checkpoint sequencing~\cite{lutris}. RAI's narrower application permits epoch-level agreement, but that choice alone is not novel. Its contribution is the report-based preservation argument required when there is no continuous certificate sequence to drain.

\para{Contributions.}
First, RAI specifies and establishes a report-exposure invariant for latent and late account finality. Second, it gives a deterministic, evidence-committed checkpoint construction that separates recovery locks from a narrow checkpoint-finality rule for uniquely retained notarized prefixes and reconciles checkpoint state with live finality. Third, it composes that construction with lagged committees and a restricted cross-epoch overlap rule. The closing committee decides one valid checkpoint; the already known successor independently verifies and installs it. An existing Kudzu-based voting construction supplies agreement~\cite{kudzu}; its mechanics and the single-writer restriction are not claimed as new contributions.

\para{Scope of the claims.}
The main protocol uses a conservative overlap rule that covers every unresolved position being finalized, including the certificate's target. A fresh child does not inherit cross-epoch safety merely because its parent was notarized twice. Supplement~B separately specifies the temporal premise of a more aggressive predecessor-backed optimization; that optimization is not needed by the main safety theorem. The results are protocol arguments, not machine-checked proofs. Section~\ref{sec:eval} reports single runs of a prototype on one host; they illustrate behavior under load but do not validate the proofs' assumptions.
